\chapter{Our Instruction Set: rvv\_ext}\label{r04:ch:rvv-ext}

\begin{r04fork}[Ours: only on the branch \texttt{rvv-support-legacy} of the pystencils submodule]
  \file{src/pystencils/backends/riscv_ext_instruction_sets.py} (248 lines),
  \file{src/pystencils/include/riscv_v_ext_prefetch.h},
  \file{riscv_v_ext_vls128.h}, \file{riscv_v_ext_vls256.h},
  \file{tests/test_rvv_qemu.py}, and three added lines in
  \file{src/pystencils/backends/simd_instruction_sets.py}.
\end{r04fork}

\isa{rvv_ext} is upstream's \isa{rvv} with five optional extras. This chapter explains
the rules it follows, each option, and how the options are implemented without changing
the vectorizer or the printer of \cref{r04:ch:vectorizer}.

\section{Rules}\label{r04:sec:ext-rules}

\begin{enumerate}
  \item \textbf{Next to upstream, not instead of it.} \isa{rvv} stays byte for byte as
    released. \isa{rvv_ext} is a new module, registered with one branch in the
    dispatcher:
\begin{pycode}
    elif instruction_set in ['rvv']:
        return get_vector_instruction_set_riscv(type_name, instruction_set)
    elif instruction_set.startswith('rvv_ext'):                              # added
        return get_vector_instruction_set_riscv_ext(type_name, instruction_set)  # added
\end{pycode}
  \item \textbf{Without options, exactly upstream's code.} \isa{rvv_ext} copies
    upstream's dictionary builder (and imports its helper \code{get_argument_string}).
    A test checks that both print identical kernels, so every difference measured
    between \isa{rvv} and an \isa{rvv_ext} variant comes from the options in its name.
  \item \textbf{Options in the name.} The options are appended to the name with
    \code{-}, the way \ps{}~1.4 names a fixed-width SVE as \isa{sve256}. A name passes
    unchanged through every layer that carries an instruction set: waLBerla's
    \code{cpu_vectorize_info}, \code{CODEGEN_CFG}, CMake lists, file and target names.
  \item \textbf{No change to the vectorizer or the printer.} Everything is done in the
    dictionary.
\end{enumerate}

\section{Names}

\code{parse_options} turns a name into options and rejects anything it does not know:
\begin{center}\small
\begin{tabularx}{\linewidth}{@{}l>{\raggedright\arraybackslash}X@{}}
  \toprule
  \textbf{token} & \textbf{meaning}\\
  \midrule
  \isa{vls128}, \isa{vls256} & vector-length specific code for VLEN 128 or 256
    (\cref{r04:sec:ext-vls}); default: vector-length agnostic\\
  \isa{m2}, \isa{m4} & LMUL 2 or 4 (\cref{r04:sec:ext-lmul}); default \isa{m1}\\
  \isa{vf} & scalar operands in \code{.vf} forms (\cref{r04:sec:ext-vf-fma})\\
  \isa{fma} & fused multiply-add (\cref{r04:sec:ext-vf-fma})\\
  \isa{pf<N>} & software prefetch $N$ bytes ahead (\cref{r04:sec:ext-pf})\\
  \bottomrule
\end{tabularx}
\end{center}
\begin{pycode}
>>> parse_options('rvv_ext-vls256-vf-fma-pf128')
{'vlen': 256, 'lmul': 1, 'vf': True, 'fma': True, 'prefetch': 128}
>>> parse_options('rvv_ext-avx')
ValueError: unknown option 'avx' in instruction set rvv_ext-avx
\end{pycode}
The order of the tokens does not matter to \code{parse_options}. The benchmark scripts
always write them in the order of the table, and check before configuring that every
name they generate means to \ps{} what they intend (\cref{r04:sec:scripts}).

\section{\texttt{vls128}, \texttt{vls256}: fixed vector length}\label{r04:sec:ext-vls}

With \isa{vls<V>} the dictionary's width is a number, $V\cdot\text{LMUL}/\text{SEW}$
(8 floats or 4 doubles at VLEN 256 and LMUL 1), and the active length of every
intrinsic is that same literal. The store template then has no \code{{loop_stop}}, so
the vectorizer cuts the loop into a vector part and a scalar remainder by itself, as it
does for AVX (\cref{r04:sec:loop}). \code{y = 0.5*x + z} becomes:
\begin{ccode}
for (int64_t ctr_0 = 0; ctr_0 < (int64_t)((_size_x_0) / (8)) * (8); ctr_0 += 8)
{
   __riscv_vse32_v_f32m1(&_data_y[ctr_0], __riscv_vfadd_vv_f32m1(__riscv_vfmul_vv_f32m1(
      __riscv_vfmv_v_f_f32m1(0.5f, 8), __riscv_vle32_v_f32m1(& _data_x[ctr_0], 8), 8),
      __riscv_vle32_v_f32m1(& _data_z[ctr_0], 8), 8), 8);
}
for (int64_t ctr_0 = (int64_t)((_size_x_0) / (8)) * (8); ctr_0 < _size_x_0; ctr_0 += 1)
{
   _data_y[ctr_0] = 0.5f*_data_x[ctr_0] + _data_z[ctr_0];
}
\end{ccode}
The active length is now a constant, so the compiler can set \code{vl} once instead of
per iteration, and the loop does not depend on the hardware's choice for the last two
vectors (\cref{r04:sec:avl}). The price: the code is only correct on hardware with at
least that VLEN. The header \file{riscv_v_ext_vls256.h}, which every such kernel
includes, makes this a compile error rather than silent wrong results:
\begin{ccode}
#if !defined(__riscv_v_min_vlen) || __riscv_v_min_vlen < 256
#error "rvv_ext-vls256 kernels need -march=rv64gcv_zvl256b (or a larger zvl)"
#endif
\end{ccode}
This is the same contract as the 2.x backend (report~1, \cref{r01:ch:rvv-backend}),
which is vector-length specific always.

\section{\texttt{m2}, \texttt{m4}: register groups}\label{r04:sec:ext-lmul}

With \isa{m2} or \isa{m4} every type and intrinsic suffix uses that LMUL
(\code{vfloat32m2_t}, \code{__riscv_vfadd_vv_f32m2}, \code{__riscv_vsetvlmax_e32m2}),
and the mask type follows the rule \code{vbool}$N$\code{_t} with $N = \text{SEW}/\text{LMUL}$
(\code{vbool16_t} for float at \isa{m2}). One instruction then processes two or four
registers' worth of elements, but only 16 or 8 vector registers remain. Large kernels
spill; in 2.x LMUL above 1 was slower on the K1. The option exists for completeness and
is tested, but no benchmark variant uses it.

\section{\texttt{vf} and \texttt{fma}: decisions at print time}\label{r04:sec:ext-vf-fma}

\paragraph{What they do.} \isa{vf} passes an operand that is the same in every lane as
a scalar to the \code{.vf} form of the operation, instead of broadcasting it:
\code{vfmul_vf(x, 0.5f, vl)} instead of \code{vfmul_vv(vfmv_v_f(0.5f, vl), x, vl)}.
That saves the broadcast (4 cycles on the X60) and a vector register. For
\code{c - v} and \code{c / v} the reversed forms \code{vfrsub_vf} and \code{vfrdiv_vf}
are used. \isa{fma} turns \code{a*b + c}, \code{c - a*b} and \code{a*b - c} into one
instruction each, \code{vfmacc}, \code{vfnmsac} and \code{vfmsac}, in the
\code{_vv} or \code{_vf} form. A fused result is rounded once instead of twice, so
it can differ from the unfused one in the last bit.

\paragraph{How.} The printer builds every vector expression by calling \code{format} on
a template, bottom up, with already printed operands (\cref{r04:sec:printer}). With
\isa{vf} or \isa{fma}, the templates for \code{'makeVecConst'} and for the four
operators are not plain strings but small \code{str} subclasses that share two tables
(abridged):
\begin{pycode}
class _Splat(_Template):          # the 'makeVecConst' template
    def format(self, *args, **kwargs):
        printed = super().format(*args, **kwargs)
        self.splats[printed] = args[0]          # remember: this C text broadcasts args[0]
        return printed

class _Operator(_Template):       # the '+', '-', '*', '/' templates
    def format(self, a, b, **kwargs):
        if 'fma' in self.forms and self.op in '+-':
            fused = self._fused(a, b, **kwargs)     # is an operand a product we printed?
            if fused is not None:
                return fused
        if 'vf' in self.forms:
            if b in self.splats:                    # vector op scalar
                printed = self.forms['vf'][self.op].format(a, self.splats[b], **kwargs)
            elif a in self.splats:                  # scalar op vector: reversed form
                printed = self.forms['vf_reversed'][self.op].format(b, self.splats[a], **kwargs)
        ...                                         # else: the plain vv template
        if self.op == '*':
            self.products[printed] = factors        # remember: this C text is x*y
        return printed
\end{pycode}
For \code{y = 0.5*x + z} the printer first prints the broadcast of 0.5;
\code{_Splat} records that the C text \code{__riscv_vfmv_v_f_f32m1(0.5f, ...)} stands
for the scalar \code{0.5f}. Then it prints the product with that text as one operand;
\code{_Operator('*')} finds it in the splat table and prints \code{vfmul_vf(x, 0.5f)}
instead, recording the product's factors. Then it prints the sum; \code{_Operator('+')}
finds one operand in the product table and prints \code{vfmacc_vf(z, 0.5f, x)}. The
broadcast and the product text are printed but end up unused, so they never reach the
C code. The result, with the prefix \code{__riscv_}, the suffix \code{_f32m1} and the
active length left out:
\begin{plaincode}
rvv:             vse32_v(&y[i], vfadd_vv(vfmul_vv(vfmv_v_f(0.5f), vle32_v(&x[i])), vle32_v(&z[i])))
rvv_ext-vf:      vse32_v(&y[i], vfadd_vv(vfmul_vf(vle32_v(&x[i]), 0.5f), vle32_v(&z[i])))
rvv_ext-vf-fma:  vse32_v(&y[i], vfmacc_vf(vle32_v(&z[i]), 0.5f, vle32_v(&x[i])))
\end{plaincode}

\paragraph{Why this is enough.} In \ps{}~1.4 hoisted constants stay scalar declarations
and every use inside the loop is printed through \code{'makeVecConst'}, and the printer
sees products as SymPy \code{Mul} terms inside an \code{Add}. So the decision is local:
whether an operand was printed as a broadcast or as a product. No analysis of the AST is
needed. One subtlety: when the vectorizer cuts a loop (\isa{vls}), it deep-copies the
loop, and with it the dictionary's templates; the shared tables would then be split
between copies. \code{_Template} therefore returns itself from \code{__copy__} and
\code{__deepcopy__}.

\section{\texttt{pf<N>}: software prefetch}\label{r04:sec:ext-pf}

With \isa{pf<N>} every load and store template is wrapped in a C comma expression that
first prefetches the address $N$ bytes ahead:
\begin{plaincode}
loadU:   (rvv_ext_prefetch_read({0}, N), __riscv_vle32_v_f32m1({0}, vl))
storeU:  (rvv_ext_prefetch_write({0}, N), __riscv_vse32_v_f32m1({0}, {1}, vl))
\end{plaincode}
The comma operator keeps the template an expression, so the printer can use it wherever
it used the plain load. The helpers (\file{riscv_v_ext_prefetch.h}) emit the Zicbop
instructions \code{prefetch.r} and \code{prefetch.w}~\cite{riscv-cmo} in their raw
encoding:
\begin{ccode}
inline void rvv_ext_prefetch_read(const void * p, long offset)
{
   __asm__ volatile("ori x0, %0, 1" : : "r"((const char *) p + offset));   // prefetch.r 0(p+offset)
}
inline void rvv_ext_prefetch_write(const void * p, long offset)
{
   __asm__ volatile("ori x0, %0, 3" : : "r"((const char *) p + offset));   // prefetch.w 0(p+offset)
}
\end{ccode}
Zicbop encodes its prefetches as \code{ori} with destination \code{x0}, an instruction
without effect, so cores without Zicbop simply execute nothing and no \code{-march} flag
is needed. The X60 executes them as prefetches although its ISA string does not list
Zicbop (measured in \file{docs/notes3}, where software prefetch gave \rfourx{2.0} with
the 2.x backend). With $N = 128$ the prefetch runs two 64-byte cache lines ahead.

\section{What the options do to a real kernel}

Counted in the D3Q19 SRT AA odd kernel in float64 (\file{scripts/notes8/lbm_qemu_check.py}):
\begin{center}\small
\begin{tabularx}{\linewidth}{@{}l>{\raggedright\arraybackslash}X@{}}
  \toprule
  \textbf{instruction set} & \textbf{arithmetic in the loop body}\\
  \midrule
  \isa{rvv} & 144 \code{vfmul}, 125 \code{vfadd}, 1 \code{vfdiv}, 110 broadcasts\\
  \isa{rvv_ext-vf-fma} & 85 \code{vfmacc_vf}, 39 \code{vfadd_vv}, 36 \code{vfmul_vv},
    23 \code{vfmul_vf}, 2 other; no broadcasts\\
  \isa{rvv_ext-vf-fma-pf128} & the same, plus 19 read and 23 write prefetches (one per
    load and store)\\
  \bottomrule
\end{tabularx}
\end{center}
About 380 operations become about 185.

\section{Compared with the 2.x backend}

Every option exists in our 2.x backend too, but there each needed real compiler work:
\begin{center}\small
\begin{tabularx}{\linewidth}{@{}>{\raggedright\arraybackslash}p{2.4cm}>{\raggedright\arraybackslash}X>{\raggedright\arraybackslash}X@{}}
  \toprule
  \textbf{feature} & \textbf{2.x backend} & \textbf{\isa{rvv_ext} (1.4)}\\
  \midrule
  fixed VLEN & always; targets \code{RISCV_RVV256} and \code{RISCV_RVV128} & option
    \isa{vls<V>}; default vector-length agnostic\\
  LMUL & option \code{cpu.vectorize.rvv_lmul} & option \isa{m2}, \isa{m4}\\
  \code{.vf} operands & always on in \code{SelectIntrinsicsRvv}: hoisted constants are
    vector symbols there, so the selector tracks which are uniform, declares scalar twins
    and checks for reassignment (about 100 lines) & option \isa{vf}: a table lookup at
    print time\\
  FMA & always on: products recognized in already lowered intrinsic calls & option
    \isa{fma}: the printer sees the SymPy product directly\\
  prefetch & option \code{cpu.vectorize.prefetch_distance}: a pass
    (\code{InsertPrefetches}, 118 lines) and platform hooks; one prefetch per distinct
    access & option \isa{pf<N>}: wrapped load and store templates; one prefetch per load
    and store\\
  row pointers & \code{HoistRowPointers}: one base pointer per access including its $x$
    offset & not ported; 1.4 already hoists one base pointer per field and row\\
  \code{in_register} & constants kept in floating-point registers instead of reloaded
    from memory & not ported\\
  \bottomrule
\end{tabularx}
\end{center}
Row pointers and \code{in_register} gave $+2.5$\,\% together in 2.x. Whether they are
worth porting is to be measured on the legacy code first, since 1.4 already does part
of what row pointers do.
